\section{First passage as an exact transport map}\label{sec:transport}

Let $D$ be a nonempty subset of $\N_+$ and $F:D\to D$ any map.
For a real $x\geq1$, define the first entrance time into $D\cap[1,x]$ by
\[
 t_x(n)=\inf\{j\in\N_0:F^j(n)\leq x\},\qquad\inf\varnothing=+\infty.
\]
Adjoin a point $\dd\notin D$, and define $p_x:D\cup\{\dd\}\to
D\cup\{\dd\}$ by
\[
 p_x(n)=\begin{cases}F^{t_x(n)}(n)&t_x(n)<\infty,\\
 \dd&t_x(n)=\infty,\end{cases}\qquad p_x(\dd)=\dd.
\]
The absorbing point $\dd$ records trajectories that never enter
$D\cap[1,x]$.  The image of $p_x$ is contained in
$(D\cap[1,x])\cup\{\dd\}$, and every point of this set is fixed by $p_x$.
The fibre $p_x^{-1}(\dd)\cap D$ consists exactly of their starting values
and may be empty.

\begin{lemma}[Nested passage]\label{lem:nested}
For $1\leq x\leq y$,
\begin{equation}\label{eq:nested}
 p_x\circ p_y=p_x=p_y\circ p_x.
\end{equation}
When $t_x(n)<\infty$, both hitting times are finite and
\[
 t_x(n)=t_y(n)+t_x(p_y(n)),\qquad
 F^{k}(p_x(n))=F^{k+t_x(p_y(n))}(p_y(n))\quad(k\geq0).
\]
\end{lemma}
\begin{proof}
If the orbit never reaches $y$, it cannot reach $x$, and the first
identity gives $\dd$ on both sides.  If it reaches $y$, the orbit prior
to that first visit is above $y$, hence above $x$.  Searching for the
first visit to $x$ can therefore begin at $p_y(n)$, giving the hitting-time
identity when finite and $\dd$ on both sides otherwise.  The equality
$p_y p_x=p_x$ follows because a successful output of $p_x$ is at most
$x\leq y$, and $\dd$ is fixed.  The identity of indexed tails follows
from the semigroup law for $F$.
\end{proof}

For a signed summable measure $\lambda$ on $D\cup\{\dd\}$, define its
pushforward $K_x\lambda=(p_x)_*\lambda$, explicitly
\[
 (K_x\lambda)(z)=\sum_{n:p_x(n)=z}\lambda(n).
\]
These maps preserve total mass, preserve nonnegativity, and contract
$\ell^1$: $\normone{K_x\lambda}\leq\normone{\lambda}$.  Indeed, take
absolute values inside each fibre sum and sum over the disjoint fibres.
Lemma~\ref{lem:nested} gives
\begin{equation}\label{eq:kernel}
 K_xK_y=K_x=K_yK_x\quad(x\leq y),\qquad K_x^2=K_x.
\end{equation}
The kernel of $K_x$ consists exactly of signed measures whose sum on
each fibre of $p_x$ is zero.  In particular,
$K_x(\delta_n-\delta_{n'})=0$ whenever $p_x(n)=p_x(n')$.

\begin{proposition}[Finite-chain transport bound]\label{prop:chain}
Let $1\leq x_0\leq\cdots\leq x_J$ and let $\nu_0,\ldots,\nu_J$
be probability measures satisfying $K_{x_i}\nu_i=\nu_i$.
Put $e_i=\normone{\nu_i-K_{x_i}\nu_{i+1}}$ for $0\leq i<J$.
Then
\begin{equation}\label{eq:chain}
 \normone{\nu_0-K_{x_0}\nu_J}\leq\sum_{i=0}^{J-1}e_i.
\end{equation}
\end{proposition}
\begin{proof}
For $J=0$ both sides are zero.  For $J\geq1$, the exact telescoping
identity is
\[
 \nu_0-K_{x_0}\nu_J=
 \sum_{i=0}^{J-1}K_{x_0}(\nu_i-K_{x_i}\nu_{i+1}).
\]
Here $K_{x_0}\nu_0=\nu_0$ and $K_{x_0}K_{x_i}=K_{x_0}$ by
\eqref{eq:kernel}.  The triangle inequality and contraction prove
\eqref{eq:chain}.
\end{proof}

The following scale formulation records exactly the summability used in
Tao's reduction.  Fix $\alpha>1$.  Suppose $\mu_s$, for $s\geq1$, is a probability
measure on $D\cap[s,s^\alpha]$ whenever that set is nonempty.  Choose
$s_*\geq2$ so every required block is nonempty for $s\geq s_*$.
For such $s$, set
\[
 \nu_s=K_s\mu_{s^\alpha},\qquad d(s)=\nu_s(\{\dd\}),\qquad
 e(s)=\normone{\nu_s-K_s\nu_{s^\alpha}}.
\]
For $R\geq s_*$ define the possibly infinite quantity
\begin{equation}\label{eq:error-tail}
 \mathcal E(R)=\sup_{u\geq R}
 \left(d(u)+\sum_{i=0}^{\infty}e(u^{\alpha^i})\right).
\end{equation}
Here $d(s)$ is the failure mass at scale $s$, while $e(s)$ measures the
change in the entrance law between two successive starting scales.

\begin{theorem}[Summable transport implies tight orbit minima]
\label{thm:transport}
For $K\geq s_*^{\alpha^3}$ and any nonempty block supporting $\mu_s$,
\begin{equation}\label{eq:abstract-tight}
 \mu_s\{n:F_{\min}(n)>K\}\leq\mathcal E(K^{1/\alpha^3}).
\end{equation}
In particular, if $\mathcal E(R)\to0$ as $R\to\infty$, the laws of
$F_{\min}$ under these block measures are uniformly tight.
\end{theorem}
\begin{proof}
Let $A_K=\{n\in D:F_{\min}(n)\leq K\}$, so $\dd\notin A_K$.
A successful passage landing in $A_K$ certifies that the original
orbit enters $A_K$.  Thus $p_u^{-1}(A_K)\subseteq A_K$ on the extended
space, for every $u$.

If $s^\alpha\leq K$, the left side of \eqref{eq:abstract-tight} is zero.
Otherwise choose $J\geq1$ minimal such that
$y=s^{\alpha^{-J}}<K^{1/\alpha}$.  Then
\[
 K^{1/\alpha^2}\leq y<K^{1/\alpha},\qquad
 u=y^{1/\alpha}\geq K^{1/\alpha^3}\geq s_*,\qquad
 u^{\alpha^J}=s^{1/\alpha}.
\]
The successful part of $\nu_u$ is supported below $u<K$, so
$\nu_u(A_K)=1-d(u)$.
For each intermediate scale $v=u^{\alpha^i}$,
\[
 \nu_{v^\alpha}(A_K)
 \geq (K_v\nu_{v^\alpha})(A_K)
 \geq \nu_v(A_K)-e(v).
\]
Summing for $0\leq i<J$ gives
\[
 \nu_{s^{1/\alpha}}(A_K)
 \geq1-d(u)-\sum_{i=0}^{J-1}e(u^{\alpha^i}).
\]
But $\nu_{s^{1/\alpha}}=K_{s^{1/\alpha}}\mu_s$, and
$p_{s^{1/\alpha}}^{-1}(A_K)\subseteq A_K$, so
$\mu_s(A_K)$ is at least this large.  Formula \eqref{eq:error-tail}
proves the bound.
\end{proof}

For example, if $d(s)\leq Ds^{-b}$ and $e(s)\leq E(\log s)^{-c}$,
with $D,E,b,c>0$, then
\begin{equation}\label{eq:power-error}
 \mathcal E(R)\leq DR^{-b}
 +\frac{E}{1-\alpha^{-c}}(\log R)^{-c}.
\end{equation}
This follows by summing the geometric series $\sum_{i\geq0}\alpha^{-ci}$.
The same theorem also accepts, for instance,
$e(s)=O((\log\log s)^{-1-c})$ for large $s$: for sufficiently large
$R>e$, writing
$h=\log\alpha>0$ and $v=\log\log R$, the sum is bounded by
$v^{-1-c}+v^{-c}/(ch)$ by integral comparison.  Tao already notes this
weaker sufficient error regime after Proposition~1.11 \cite{Tao7}.
